\firstslide{Caml}

\vspace*{2cm}

\begin{center}

{\large \bf \em Pierre Weis}

\vspace*{3cm}

Lyon -- May $22$ 2002 

\end{center}

\newslide{Plan}

\begin{itemize}
\item Past
\item Present
\item Future
\end{itemize}


\newslide{History}

1978: design of ML, first interpretor (Robin Milner)\\
1981: first ML compiler at INRIA (Guy Cousineau, G�rard Huet)\\
1984: pattern matching (Guy Cousineau)\\
1984: first class references (Pierre Weis).\\
1985: first V1.17 Caml compiler (Ascander Suarez)\\
1990: Caml V3.1, ultimate version based on JIT (Pierre Weis)\\


\newslide{Caml V3.1}

- general ``or'' patterns (binding variables), \\
- general ``alias'' patterns (pat as pat),  \\
- guarded patterns (when construct), \\
- lazy values (for sum and product types) with
automatic forcing at access time during pattern matching, \\
- declared mutable constructors and labels, with proper contents
binding during pattern matching, \\
- dynamic values with coercion via pattern matching, \\
- parsing facilities included in the language (grammars), \\
- pretty-printing facilities (printf and format).


\newslide{Caml V3.1 extensions}

- extensible sum types (new constructors can be added after
the initial definition of the type), \\
- polymorphic overloading with static resolution
(overloading accepted iff there is only one acceptable solution in the
context), \\
- fully typechecked forward declarations for inter-modules
recursion (polymorphic forwards used type scheme unification), \\
- introspection facilities (lexical and syntactic macro generation), \\
- generalized looping construct, \\
- a polymorphic ``print'' primitive.

\newslide{Past of Objective Caml}

1991: Caml Light compiler on PC (Xavier Leroy, Damien Doligez)\\
1996: modules, functors, native code compiler (Xavier Leroy)\\
1998: objects $\rightarrow$ Objective Caml (J�r�me Vouillon)\\
2000: labels, optional arguments, polymorphic variants (Jacques Garrigue)\\
2002: dynamically linked libraries, libraries encapsulation (Xavier Leroy)


\newslide{Present}

{\bf Compiler}: Objective Caml has a fast byte-code {\em and} native
compilers (both Caml Light and Caml inheritance). Compiles quickly
programs that go fast!

{\bf Runtime}: Objective Caml has a state of the art GC
(generational, compacting, ``malloc'' compatible, static allocation,
finalized objects, and weak pointers).

{\bf Library}: Objective Caml has a mature user's library polished by
20 years of research. 


\newslide{Programming in the small}

{\bf Printf}: statically typed printing facility with compact formats.

{\bf Toplevel}: interactive system that allows the programmer to
load and execute pieces of code.

{\bf Test}: the toplevel allows testing and simple debugging (trace).

Also convenient for teaching \ldots


\newslide{Programming in the large}

Objective Caml has a powerful module system that statically verifies
modules compatibility.

Other features for programming in the large:
\begin{citemize}
\item Value I/Os: any value can be written to a file that can be read
back (architecture aware input/output)
\item Camllex Camlyacc
\item Architecture independant graphics library
\item GUIs
\end{citemize}

\newslide{Other features}

{\bf Debugging}: Objective Caml has a {\em replay} debugger.

{\bf FI}: Objective Caml has a foreign interface with C and Fortran.

{\bf Portability}: extremely high (compilers are ported and easy
to port, users' programs are assumed to be portable at no cost)

{\bf Users' communauty}: Objective Caml has a reactive and
active users' communauty all over the world. (Not millions, but
sophisticated programmers.)

\newslide{Strength of Objective Caml}

{\bf State of the art in language design}: ``The language of choice for the
discriminating hackers!''

{\bf ICFP programming contest}:\\
1999: 2nd (after a specialized C program)\\
2000: 1st and 2nd\\
2001: 2nd (after a Haskell program) \\
2002: 1st

{\bf Benchmarks}: consistantly faster than any other functional programming
languages, between C++ and C-O2

\newslide{Strength of Objective Caml}

{\bf Static checkings}.

\begin{citemize}
\item {\bf Type reconstruction} (no annotations in programs)
\item {\bf Genericity} (structural polymorphism)
\item {\bf Easy} value definition and manipulation
\item {\bf Theorems} due to clean semantics
\end{citemize}

\newslide{Type reconstruction}

No type annotation in programs,\\
Type modifications are automatically propagated.

{\bf Theorem}: no type error can occur at runtime.

\begin{small}\begin{alltt}
# let succ x = x + 1;;
val succ : int -> int = <fun>
# let apply1 x = x 1;;
val apply1 : (int -> 'a) -> 'a = <fun>
\end{alltt}\end{small}

{\bf Theorem}: the type (scheme) found by the compiler is the most general.\\

\newslide{Type reconstruction}

Simplifying the definition.

\begin{small}\begin{alltt}
# let const x = 1;; \pause
val const : 'a -> int = <fun> \pause
# let id x = x;; \pause
val id : 'a -> 'a = <fun> \pause
# let f x = x ( + );; \pause
val f : ((int -> int -> int) -> 'a) -> 'a = <fun> \pause
# let g x = ( + ) 1;; \pause
val g : 'a -> int -> int = <fun> \pause
# let h = ( + ) 1;; \pause
val h : int -> int = <fun> \pause
# let i = ( + );; \pause
val i : int -> int -> int = <fun>
\end{alltt}\end{small}

\newslide{Type reconstruction}

Working with functions

\begin{small}\begin{alltt}
# let compose f g = function x -> f (g (x));; \pause
val compose : ('a -> 'b) -> ('c -> 'a) -> 'c -> 'b = <fun> \pause
# compose succ (function s -> String.length (s ^ s));; \pause
- : string -> int = <fun> \pause
# compose (function s -> String.length (s ^ s)) succ;; \pause
                                                ^^^^
This expression has type int -> int but is here used with type int -> string \pause
# let z x = (compose compose compose) x;; \pause
val z : ('a -> 'b) -> ('c -> 'd -> 'a) -> 'c -> 'd -> 'b = <fun> \pause
\end{alltt}\end{small}

No problemo: it works!

\newslide{Easy value definition}

{\bf Powerful} data type definitions.

{\bf Implicit} allocation and deallocation.

{\bf Compact} syntax for creation.

{\bf Theorem}: no value can be left uninitialized.

\newslide{Type definitions}

{\bf Product} types: cartesian products (\verb"t1 * t2") or {\bf records}
\begin{alltt}
type monom = \{coefficient : float; degree : int\}
\end{alltt}

{\bf Sum} types:
\begin{alltt}
type expression =
   | Const of float
   | Sum of expression * expression
\end{alltt}

\newslide{Type definitions}

Type {\bf abbreviations}: {\tt type polynom = monom list}

{\bf ADT}: defined in the implemention of the module, but
definition is not exported in the interface

Any kind of type can be polymorphic and/or recursive
\begin{alltt}
type 'a polylist = Nil | Cons of 'a * 'a polylist
\end{alltt}

\newslide{Value creation}

\begin{small}\begin{alltt}
# let m1 = \{coefficient = 3.14; degree = 2\};;
val m1 : monom = \{coefficient = 3.14; degree = 2\}
\end{alltt}

\begin{alltt}
# let e1 = Sum (Const 1.0, Const 2.0);;
val e1 : expr = Sum (Const 1., Const 2.)
\end{alltt}

\begin{alltt}
# let l1 = 1 :: 2 :: [];;
val l1 : int list = [1; 2]
# let rec linfty = 1 :: linfty;;
val linfty : int list =
  [1; 1; 1; 1; 1; 1; 1; 1; 1; 1; 1; 1; 1; 1; 1; 1; 1; 1; 1;
   ...]
\end{alltt}
\end{small}
{\bf Theorem}: once created, no value can be unadvertedly modified.

\newslide{Easy value manipulation}

Pattern matching
\begin{small}\begin{alltt}
let rec add_pol (p1, p2) =\pause
  match p1, p2 with\pause
  | [], p2 ->\pause p2\pause
  | p1, [] -> p1\pause
  | (m1 :: r1), (m2 :: r2)\pause when m1.degree < m2.degree ->\pause
      m1 :: add_pol (r1, p2)\pause
  | (m1 :: r1), (m2 :: r2) when m1.degree = m2.degree ->\pause
      \{coefficient = m1.coefficient +. m2.coefficient;\pause
       degree = m1.degree\} :: add_pol (r1, r2)\pause
  | _, _ ->\pause add_pol (p2, p1);;\pause
\end{alltt}\end{small}

\newslide{Structural polymorphism}

Follows values structure
\begin{small}\begin{alltt}
# [];;\pause
- : 'a list = []\pause
# [] :: [];;\pause
- : 'a list list = [[]]\pause
\end{alltt}\end{small}

\begin{small}\begin{alltt}
let rec sort l = \pause
  match l with\pause
  | [] -> []\pause
  | e :: l -> insert e (sort l)\pause
and insert e l =\pause
  match l with\pause
  | [] -> [e]\pause
  | x :: l -> if e <= x then e :: l else x :: insert e l;;\pause
val sort : 'a list -> 'a list = <fun>
val insert : 'a -> 'a list -> 'a list = <fun>
\end{alltt}\end{small}

\newslide{Exceptions}

Exists in Caml from the very beginning.

-- Fully type-checked.\\
-- No counterparts in function types.\\
-- Raised using {\tt raise}, caught with {\tt try}.\\
-- Selection via pattern matching.
\begin{small}\begin{alltt}
try
  let oc = open_out "/tmp/foo" in
  ...
with
| System_error s -> eprintf ("System error " ^ s)
| x -> eprintf "Spurious exception"; raise x
\end{alltt}\end{small}
In addition, exceptions are first class citizens!

\newslide{Code safety}

{\bf Static verifications}.

{\bf Strong typing} is all over the place (inside modules, between modules).

{\bf Code structuration}: data type definitions are mandatory.

{\bf Static binding} (even {\em hyper-static} binding).


\newslide{Code safety}

{\bf Theorem}: once defined, no binding can be modified.

{\bf Corollary}: hypotheses that hold when writing a function, still
hold when executing it.

Enables {\em true} functions (so-called closures) as first class citizens.

\newslide{Code evolutivity}

{\bf Pattern matching} is a concise ``test and bind'' facility.

Its clean and precise semantics leads to extra static sanity checks:\\
- {\bf exhaustivity checking}\\
- {\bf unused match case detection}

Efficient way to consistantly propagate modifications into the code:
modify and compile!

\newslide{Code evolutivity}

Multiparadigmatic language: you're not stuck to a given style.

\pause Smooth integration of
\begin{center}
\pause\textcolor{blue}{\advitransbox{slide,from=right}{imperative,}}
\pause\textcolor{green}{\advitransbox{slide,from=left}{functional, and}}
\pause\textcolor{red}{\advitransbox{slide,from=top}{object}}
\pause\textcolor{black}{\advitransbox{slide,from=bottom}{paradigms.}}
%imperative, functional, and objects paradigms.
\end{center}

\pause
On top: {\advitransbox{block,steps=50}{a first order functional language}}\\
\pause{\advitransbox{wipe,steps=150,from=left}{for powerful module programming}}

\pause(\textcolor{red}{%
\advitransbox{wipe,steps=100,from=right}{\em templates on steroids}} ?).

\newslide{Strength of Objective Caml}

The compiler is good at {\bf verifying program coherence}.

Usual programmers' experience with Caml:

\begin{center}
{\em Once the compiler ``accepts'' to compile the program,\\
 the program works!}
\end{center}

\pause Unfortunately wrong, but \ldots almost true!

\newslide{Weaknesses}

{\bf Technically}:~\pause None that I know of :)\pause\\
{\bf Socially}:\pause
\begin{citemize}
\item Not like SML (standard ?)\pause
\item Not like C++ (large and uneasy to understand ?)\pause
\item Not enough advertized (lacks the Java ``evangelists'' ?)\pause
\item Not expansive enough (as Maple and Mathlab ?)\pause
\end{citemize}
Not American enough ? %To be french ?\pause

\newslide{Strength of Objective Caml}

\textcolor{red}{\bf Clean semantics} (both static and dynamic).

\textcolor{red}{\bf Clean implementation} (both compilers, additional tools, and libraries).

\textcolor{red}{\bf Encourage} well structured and well written code.

\textcolor{red}{\bf Allows} easy prototyping (type reconstruction).

\textcolor{red}{\bf Helps} for programming in the large (modules).

\newpage\advibg{image=fumicamltrans.ppm}

\pause\stitle{Future}

%\let\orgheader\adviheader
%\def\adviheader{\orgheader\advibg{image=fumicamltrans.ppm}}

\pause Overloading and programming with types ?

\vfill

\pause {\Huge{On ne voit rien !}}

\newpage\advibg{alpha=0.3,image=fumicamltrans.ppm}

\pause\stitle{Future}

\pause Overloading and programming with types (dynamic values) ?

\pause Automatic test generation ?

\pause More static analyses in the code ?

\pause Mixins ?

\pause Long term: programming with proofs ?

\newslide{Bibliography}

``Le langage Caml'',\\
Pierre Weis et Xavier Leroy, InterEditions 1999.

``D�veloppement d'applications avec Objective Caml'',\\
Chailloux, Manoury, Pagano, O'Reilly 2000.

Web:

\verb"http://pauillac.inria.fr/caml/index-fra.html"\\
\verb"http://pauillac.inria.fr/caml/FAQ/index-fra.html"\\


\newslide{What's that ?}

\advitransition{slide}

\def\tpause{\advitransition{none}\pause}
\pause Xdvi ?%
\tpause{\advitransbox{slide,from=right}{\textcolor{red}{No}}}%
\tpause~Ghostview ?%
\tpause{\advitransbox{slide,from=right}{\textcolor{red}{No}}}%
\tpause~Acroread ?%
\tpause{\advitransbox{slide,from=right}{\textcolor{red}{No}}}%
\tpause~Powerpoint ?%
\tpause{\advitransbox{slide,from=right}{\textcolor{red}{No}}}

\tpause Its \textcolor{blue}{\bf Active-DVI !}
\tpause Entirely written in Caml, \textcolor{red}{\bf hence~}\tpause reliable and easy to extend.
\tpause Entirely programmable, this is:\tpause

\bigskip
\begin{center}
\textcolor{blue}{\Large{\bf The presenter of choice}}

\textcolor{blue}{\Large{\bf for the discriminative hacker!}}
\end{center}

\newslide{That's all, folks!!!}
\advitransition{wipe}

\begin{center}
\includegraphics[width=0.25\textwidth]{advilogo.eps}
\end{center}

\def\advimasquotte{\adviembed[width=1.56cm,height=1.824cm]%
{animate -geometry !g! -window !p advilogo.anim.gif}}

\hbox to \hsize{\pause\hfill\advimasquotte\pause\hfill\advimasquotte\pause\hfill\advimasquotte\pause\hfill\advimasquotte\pause\hfill\advimasquotte\hfill}

\vfill

\pause

\newslide{Interaction}

\adviembed[width=0.15\textwidth,height=0.04\textheight]{xterm -geometry !g -e ocaml}

\newslide{The code of this slide}

\begin{verbatim}
\newslide{Interaction}

\adviembed[width=0.15\textwidth,height=0.04\textheight]%
  {xterm -geometry !g -e ocaml}
\end{verbatim}
